% Generated by task tex-opening from dossiers/gt-opening-compile.md (sections A, B and C).
% With the corrections of dossiers/verify-gt-opening.md (the citation check), marked in the text.
\subsection{The opening and the compilation: the transcript as deployed}\label{sec:gen-opening-transcript}

\begin{sloppypar}
Short forms in this section. Specification files are under \file{doc/leanvm/body/}: \code{03:} is
\file{03-proving-primitives.tex}, \code{04:} is \file{04-committing-the-witness.tex}, \code{08:} is
\file{08-end-to-end-protocol.tex}, \code{A:} is \file{a-ring-switching.tex} and \code{B:} is
\file{b-polynomial-commitment-scheme.tex}. Rust files: \code{whir.rs}, \code{whir_config.rs},
\code{whir_induce.rs}, \code{whir_ntt_ext.rs} and \code{stack_open.rs} are under
\file{crates/pcs/src/}, and \code{pcs/merkle.rs} is \file{crates/pcs/src/merkle.rs}; \code{fs/} is
\file{crates/fiat_shamir/src/}; \code{mod.rs} is \file{crates/lean_vm/src/cpu/mod.rs}, and
\code{pcs.rs}, \code{witness.rs}, \code{leaf.rs} are under \file{crates/lean_vm/src/};
\code{primitives/src/} is \file{crates/primitives/src/}. \code{py:} is
\file{python-verifier/verifier.py}. These are at leanVM's pin \code{a386121f}. \code{bp:} is the
blueprint \file{docs/roadmap/protocol-blueprint.md} at leanerVM \code{b435631} (the checkout
\code{144c5aa} shifts its lines after line 168 by four); Lean files named without a directory are
under \file{LeanerVM/Protocol/} at \code{b435631}; ArkLib is cited at its old pin \code{dca90385}
unless a revision is named. A bare line number (\code{:1012}) is in the file the column heading or
the sentence names. The number after the colon is the line.
\end{sloppypar}

\begin{sloppypar}
The leanVM working tree was moved off the pin on 2026-09-30 at 09:42, not by this review (its
reflog: \code{pull: Fast-forward}, then \code{checkout: moving from main to
feat/drop-prover-public-input-challenge}; HEAD \code{248da071}). Every line cited here was read on
2026-09-29 while HEAD was \code{a386121f}, and the ones quoted were re-verified on 2026-09-30 with
\code{git show a386121f:<path>}.
\end{sloppypar}

$P→V$ is a prover message (a read of the stream, bound into the Fiat--Shamir state as it is read),
$V$ a verifier computation or check, $V→P$ a challenge (one squeeze of the chain per element of
$\E$).

% ---------------------------------------------------------------- A.1
\subsubsection{From a pooled claim to a weighted claim on the stack}\label{sec:gen-opening-pool}

A pooled claim is $\mle{P}_i(ζ) = c$ on one column; through the stacking selectors it is
$Σ_w W(w)·\mle{q}(w) = c$ with $W(w) = \eq((ζ, \mathit{sel}_i), w)$ (\code{04:14-18}; \code{08:97}:
``Each pooled claim is one evaluation of a column at a point, its value already supplied; via the
stacking selectors (§4.1) it becomes a weighted sum $Σ_w W_j(w)\,\mle{q}(w) = c_j$ over the stacked
witness. (The ring-switched claim arrives already in this form, its weight supported on the
$\qflock$ region.)'').

{\footnotesize\setlength{\tabcolsep}{4pt}
\begin{longtable}{@{}>{\raggedright\arraybackslash}p{0.14\linewidth}>{\raggedright\arraybackslash}p{0.15\linewidth}>{\raggedright\arraybackslash}p{0.39\linewidth}>{\raggedright\arraybackslash}p{0.26\linewidth}@{}}
\caption{From a pooled claim to a weighted claim, as deployed.}\label{tab:gen-opening-pool}\\
\toprule
Step & Specification & Rust & Python \\
\midrule
\endfirsthead
\toprule
Step & Specification & Rust & Python \\
\midrule
\endhead
\bottomrule
\endfoot
A column claim becomes a slot claim
  & \code{08:97}; \code{04:14-18}
  & \code{mod.rs:790-814} \code{slot_claims}: \code{SlotClaim::Point { offset, low_point, value }}
  for a committed column, \code{SlotClaim::Strided { offset, slot, stride_log, point, value }} for a
  \code{BLAKE2S} value limb (``virtual'', routed to $\qflock$)
  & \code{py:520-522} \code{ColumnClaim.on_stack}: \code{point = layout.placements[column] .stack_point(self.point, layout.stack_log)};
  \code{Placement.stack_point} (\code{py:295-297}) inserts the selector bits around the point, the
  low bits for a strided limb \\
The weight of a point claim
  & \code{04:18} $W(w) = \eq((ζ, \mathit{sel}_i), w)$
  & \code{stack_open.rs:294-303} \code{stack_claim_eq_at}: \code{eq(low_point, x[..n]) · ∏_k (sel_k ? x_k : 1 + x_k)}
  with \code{sel = offset >> n}; \code{stack_open.rs:305-323} the strided form
  \code{eq(slot bits, x[..stride]) · eq(point, x[stride..block]) · eq(sel, x[block..])}
  & \code{py:522} \code{lambda x: eq_eval(point, x)} on the full stack point \\
The ring-switched claim's weight
  & \code{08:91}, \code{A:46} $W(u) = Φ(\eq(r, u))$ on $\qflock$, lifted by its selector
  & \code{stack_open.rs:45-48} \code{b(x) = eq(sel, x_hi) · Σ_i λ^i·MLE(rs_eq_ind_i)(x_lo) + Σ_j λ^(n_rs+j)·eq(claim_j, x)};
  verifier \code{stack_open.rs:530-546}
  & \code{py:1411-1412} \code{ringswitch = (lambda x: qflock.eq_above(x) * ringswitch_weight(x[: qflock.variables]), ringswitch_target)} \\
\end{longtable}
}

The verifier never materializes a weight: it evaluates the batched weight's extension once, at
WHIR's terminal point (\code{stack_open.rs:529-548} ``Evaluate the lifted weight once, at the
terminal sumcheck point''; \code{py:1362}, \code{py:1079}). That is what ``MLE-friendly''
(\code{03:113}) buys.

% ---------------------------------------------------------------- A.2
\subsubsection{The batching challenge \texorpdfstring{$λ$}{lambda} and the assignment of its powers}\label{sec:gen-opening-lambda}

{\footnotesize\setlength{\tabcolsep}{4pt}
\begin{longtable}{@{}>{\raggedright\arraybackslash}p{0.12\linewidth}>{\raggedright\arraybackslash}p{0.2\linewidth}>{\raggedright\arraybackslash}p{0.36\linewidth}>{\raggedright\arraybackslash}p{0.26\linewidth}@{}}
\caption{The batching challenge and its powers, as deployed.}\label{tab:gen-opening-lambda}\\
\toprule
Step & Specification & Rust & Python \\
\midrule
\endfirsthead
\toprule
Step & Specification & Rust & Python \\
\midrule
\endhead
\bottomrule
\endfoot
When $λ$ is drawn
  & \code{08:98} ``drawn after every claim value is bound''
  & \code{stack_open.rs:397-400}: after the six ring-switching map challenges
  (\code{stack_open.rs:394}), \code{let lambdas = powers(ps.sample(), ring.claims.len() + point_claims.len());}
  with the comment ``Nothing is observed first: every claim value reached the caller through a
  binding stream read''; verifier \code{stack_open.rs:518}
  & \code{py:1361} \code{scales = powers(transcript.sample(), len(claims))}, after
  \code{ring_switch} drew its six (\code{py:1343}) \\
Which claim takes which power
  & \code{08:98} ``Claim $j$ takes the weight $λ^{j−1}$, the ring-switched claim first and the
  pooled column claims after''
  & \code{stack_open.rs:401} \code{let (lambdas_rs, lambdas_pd) = lambdas.split_at(ring.claims.len());}:
  the ring-switched claims take $λ^0, …$, the point claims the rest in \code{finish_claims} order
  (\code{mod.rs:656-667}: bus, table, public input)
  & \code{py:1413} \code{[ringswitch, *(c.on_stack(layout) for c in claims)]} with
  \code{claims = [*bus.claims, *table_sumcheck_claims, *public]} (\code{py:1395, 1402}) \\
Number of claims
  & $J$ (\code{08:98})
  & $1 + 112 = 113$ (one ring-switched claim, \code{mod.rs:767}; the pool as the dossier on the
  table sumcheck, \file{gt-table-pub.md}, section A.3, gives it)
  & same \\
The batched target
  & \code{08:99} $C_λ = Σ_j λ^{j−1} c_j$
  & \code{stack_open.rs:521-527} (verifier), \code{:413-415, 257, 285} (prover)
  & \code{py:1362} \code{dot(scales, values)} \\
The batched weight
  & \code{08:99} $W_λ = Σ_j λ^{j−1} W_j$, ``the verifier evaluating $W_λ$ itself''
  & \code{stack_open.rs:531-546} \code{eval_b_at}
  & \code{py:1362} \code{lambda point: dot(scales, [weight(point) for weight in weights])} \\
\end{longtable}
}

The three sources agree on every row (the dossier on Flock and ring switching,
\file{gt-flock-ring.md}, section 6, gives the ring switching; the dossier on the table sumcheck,
\file{gt-table-pub.md}, section A.3, the pool). \code{powers} starts at $λ^0$ in both
implementations (\code{primitives/src/field/mod.rs:51-61}; \code{py:186-188}), matching $λ^{j−1}$ for
$j = 1$. (The citation check corrected the dossier's \code{whir_induce.rs:99-107}, which is
\code{power_weights}, the per-query weights inside a level; both start at $λ^0$.)

% ---------------------------------------------------------------- A.3
\subsubsection{What runs after \texorpdfstring{$λ$}{lambda}: the question settled}\label{sec:gen-opening-after-lambda}

\paragraph{Specification.} \code{08:99}: ``Fold the $J$ claims into $W_λ = Σ_j λ^{j−1} W_j$ and
$C_λ = Σ_j λ^{j−1} c_j$, and run the inner-product PCS opening on $Σ_w W_λ(w)\,\mle{q}(w) = C_λ$
(Annex B), the verifier evaluating $W_λ$ itself. This one opening discharges every pooled claim;
there is no separate reduction sumcheck.'' Annex B's \code{Open} takes $J_0 ≥ 1$ weighted claims
and its level-0 step 1 is ``(batch) V sends $λ ← \E$; the batched claim is
$(W, σ_0) := (Σ_t λ^{t−1} W_t, Σ_t λ^{t−1} c_t)$'' (\code{B:108-110}): §8.5's $λ$ \textbf{is} the
level-0 batching challenge of Annex B's opening protocol (Protocol B.6 at the pin; the dossier's
``Protocol B.1'' here is the blueprint's number, as the citation check noted,
\cref{f:opening-annex-numbers}), with $J_0 = J$. Its step 2 is the level's
sumcheck, $ℓ_0$ rounds (\code{B:111}).

\begin{sloppypar}
\paragraph{Rust.} \code{stack_open.rs:452-461} hands \code{(stack, b_stack, target)} to
\code{recursive_prover_with_basis} (``One WHIR over the full stack against the combined claim'').
Its first act:
\end{sloppypar}

\begin{srccode}
    let (mut sc_prover, start_msg) =
        sumcheck_span.in_scope(|| SumcheckProver::new(witness, b_initial, target, lane_block));
    send_msg(ps, start_msg, target);
\end{srccode}
\src{crates/pcs/src/whir.rs:1324-1326 at a386121f}

\begin{sloppypar}
\code{send_msg} transmits a round polynomial (\code{whir.rs:509-511}). Then, \code{whir.rs:1328-1334},
six times: sample $r_j$, fold, send the next round polynomial. Only after that is the first
recursive root sent (\code{whir.rs:1358} \code{ps.add_root(&wtns_1.root());}). The verifier mirrors
it: \code{whir.rs:2087-2090} \code{let mut t_r = target; let Some(mut running_quad) = recv_quad(vs, t_r)},
\code{:2100} \code{replay_fold_rounds(vs, initial_k, …)}, \code{:2104} \code{vs.next_root()}.

\paragraph{Python.} \code{py:1362} calls \code{verify_whir(transcript, stack_log, log_inv_rate, dot(scales, values), root, …)},
whose first read is \code{py:1012} \code{running_quad = transcript.sumcheck_round_poly(3, target)},
then \code{py:1017-1024} the fold rounds, then \code{py:1034} the next root.
\end{sloppypar}

\paragraph{So.} In the deployed protocol there is \textbf{no sumcheck between $λ$ and WHIR}; the
first message after $λ$ is WHIR's first round polynomial. The sumcheck rounds that reduce the
weighted claim are \textbf{WHIR's own}: $μ = 6 + Σ_i k_i + n_{\mathrm{res}}$ rounds in all
(\code{whir_config.rs:748-751} ``\code{initial_k + Σ k (L1+) + yr_log_n = log_n}''), interleaved
with the level roots, the out-of-domain samples, the grinding nonces, the query positions and the
per-level batching challenges. The running claim is not $Σ W_λ·q$ throughout: at every level the
out-of-domain claim and the query-consistency claim are glued into it with powers of that level's
$λ_i$ (\code{whir.rs:1131-1154} \code{glue_pending}, verifier \code{:1681-1701}
\code{batch_level_claims}; \code{py:1058-1063}), and the protocol closes on one equation
$\mathit{weight} · \mle{f}_{\mathrm{final}}(ρ_{\mathrm{tail}}) = t_r$ (\code{whir.rs:2308};
\code{py:1082-1083}), where $\mathit{weight}$ is $\mle{W}_λ$ at the rotated terminal point plus
every glued weight (\code{whir.rs:2264-2307}). There is never an evaluation of $\mle{q}$ at a
point; $q$ is folded, and the folds are committed.

% ---------------------------------------------------------------- A.4
\subsubsection{What the blueprint specifies instead}\label{sec:gen-opening-blueprint}

\begin{sloppypar}
\begin{itemize}
\item Layer 10 (\code{bp:1112-1114}): \code{openingPhase (μ) (J) … (pSpec := V_to_P : E ; sumcheck (W_λ · q) rounds ; the final evaluation query)},
  error ``$(J − 1)/\abs{\E}$ on $λ$, $2/\abs{\E}$ per round''. The section is titled ``the claim
  pool, the opening sumcheck, and the oracle protocol'' (\code{bp:1101}), and the scope says ``the
  claim pool and the opening sumcheck, as ArkLib reductions'' (\code{bp:127}).
\item Layer 11 (\code{bp:1173}): \code{whirOpen (params) : OracleReduction … (StmtIn := Fin J → WeightedClaim μ) …},
  ``stated exactly as Protocol B.1: batch, $ℓ_i$ sumcheck rounds, commit, …'' (\code{bp:1185}).
\item Layer 12 (\code{bp:1214-1215}): ``The compiled oracle protocol: WHIR in place of the evaluation
  oracle. \code{def leanVmIopp … (OStmtIn := codeword oracles)}''.
\item The convention \emph{The oracle} (\code{bp:317}): ``\code{OracleInterface} query
  \code{Vector E μ_stack} and answer \code{eval₂Mle q}''; Layer 0 (\code{bp:637-642}) and the built
  instance (\file{LeanerVM/Protocol/Field.lean:102-106} at \code{b435631}).
\end{itemize}
\end{sloppypar}

Read together, the compiled verifier of the blueprint runs $λ$, a $μ$-round sumcheck on $W_λ·q$
(two coefficients a round), receives a value $v = \mle{q}(ρ)$, checks $\mle{W}_λ(ρ)·v$ against the
running claim, then runs \lean{whirOpen} on the one claim $(\eq(ρ, ·), v)$, which, ``stated exactly
as Protocol B.1'', starts with its own level-0 batching challenge ($J_0 = 1$) and its own
$μ$-round sumcheck. That transcript has one extra challenge, $μ$ extra round polynomials and one
extra scalar, and is not leanVM's: a \lean{verify} that is ``the Fiat--Shamir compilation of that
oracle verifier'' (\code{bp:27-29}) cannot accept a Rust proof, and \lean{verify_iff_compiled}
(\code{bp:1218-1221}) is false for every proof the pinned prover emits.

This is the divergence between the pinned §8.5 (\code{08:94-101}) and the blueprint's Layer 10
that the review set out to classify. It is \textbf{an error of the blueprint}, introduced by the
choice of the evaluation interface at Layer 0: with an evaluation oracle, the only way to
discharge a weighted claim in the ideal model is to reduce it to an evaluation, and the extra
sumcheck follows (\cref{f:opening-extra-sumcheck}).

% ---------------------------------------------------------------- A.5
\subsubsection{The alternatives, and the recommendation on the oracle interface}\label{sec:gen-opening-alternatives}

\paragraph{(i) Keep the evaluation oracle and the extra sumcheck; owe a theorem relating the two
protocols.} The ideal-model theorems stay as they are, but \lean{verify} can no longer be the
Fiat--Shamir compilation of the oracle verifier (\cref{sec:gen-opening-blueprint}). The theorem
owed would be: ``the compiled protocol \emph{with} the extra sumcheck and WHIR on an evaluation
claim is secure $⇒$ the deployed protocol is secure''. No such implication exists between two
protocols with different transcripts; each must be proved on its own, so the security of
\lean{verify} would rest on a WHIR theorem for weighted claims anyway (Theorem B.7 is stated for
$J_0$ weighted claims), and the opening sumcheck of Layer 10 would be proved and never used.
Rejected.

\begin{sloppypar}
\paragraph{(ii) The inner-product oracle interface.} A query is a weight \lean{W : Weight μ}, the
answer is $Σ_w W(w)·q(w)$: the idealization of Definition 3.13 (\code{03:112-118}: ``An
\emph{inner-product commitment scheme} commits to $g$ and later proves any such claim … An
evaluation claim $\mle{g}(r) = c$ is the case $W = \eq(r, ·)$''). The opening phase is then
``$λ$, then one weighted query at $W_λ$ compared with $C_λ$'', a one-challenge component with
error $(J − 1)/\abs{\E}$ and no sumcheck; WHIR realizes the one query. Checked at the pins (probe
\file{InnerProductOracle.lean}, \cref{sec:gen-opening-probe-inner}, exit 0 on 2026-09-29 against
ArkLib \code{dca90385}):
\begin{itemize}
\item ArkLib's \lean{OracleInterface} is a class with a query type and an answer function
  (\file{ArkLib/OracleReduction/OracleInterface.lean:55-57} at \code{dca90385}:
  \code{class OracleInterface (Message : Type u) where Query : Type v; toOC : OracleContext Query (ReaderM Message)});
  \lean{Weight μ} is a structure in \code{Type} (\file{Spine/Seams.lean:103-109} at \code{b435631}),
  so \code{Query := Weight n} is accepted and the answer is \code{W.pair q}
  (\code{Seams.lean:112-113}).
\item The evaluation query is the special case \code{eqWeight p}, by the built lemma
  \lean{eqWeight_pair} (\code{ClaimWeights.lean:51-54} at \code{b435631}): the current
  \lean{evalOracle_answer} becomes a corollary.
\item ArkLib's functional commitment scheme over this interface has for opening statement
  \code{(commitment, weight, value)} (\file{ArkLib/Commitments/Functional/Basic.lean:62-64} at
  \code{dca90385}, copied in the probe): it is an inner-product commitment scheme, which is what
  Layer 11 must instantiate (``\code{Commitment.Scheme} … The commitment interface WHIR's commit and
  open instantiate'', \code{bp:249}).
\item The opening verifier ``read $λ$, query at $W_λ$, compare with $C_λ$'' is an
  \lean{OracleVerifier} over \code{ProtocolSpec 1} with \lean{keepOracles} (the probe's
  \lean{openVerifier}).
\end{itemize}
\end{sloppypar}

Consequences, layer by layer:

{\footnotesize\setlength{\tabcolsep}{4pt}
\begin{longtable}{@{}>{\raggedright\arraybackslash}p{0.14\linewidth}>{\raggedright\arraybackslash}p{0.82\linewidth}@{}}
\caption{The inner-product oracle interface: what changes, layer by layer.}\label{tab:gen-opening-consequences}\\
\toprule
Where & Change \\
\midrule
\endfirsthead
\toprule
Where & Change \\
\midrule
\endhead
\bottomrule
\endfoot
Layer 0 (\file{Field.lean})
  & replace \lean{evalOracle} by the inner-product instance; \lean{Weight} (and \lean{Weight.pair})
  move below \file{Field.lean} (today they are in \file{Spine/Seams.lean}, above it);
  \lean{evalOracle_answer} becomes \code{answer q (eqWeight p) = eval₂Mle …}. The only consumers of
  \lean{OracleInterface.answer} on a \lean{Column} at \code{b435631} are three theorems on
  verifier-computed columns (\code{FixedColumns.lean:59-61, 97-99, 113-115}) and the tests
  (\file{tests/LeanerVMTests/Protocol/Field.lean:30-44}); the former are restated with
  \lean{eval₂Mle}, since the index and bytecode columns are never oracles. (The citation check
  found this list incomplete: \file{tests/LeanerVMTests/Protocol/FixedColumns.lean:32, 49, 76, 115}
  also use \lean{OracleInterface.answer} on columns, and the docstring of
  \file{LeanerVM/Protocol/Stack.lean:18} refers to it; the recommendation is unaffected in kind, the
  list of files to change is longer.) The convention
  \emph{The oracle} (\code{bp:317}) and Layer 0 (\code{bp:637-642}) change accordingly \\
The spine
  & nothing in the types: \lean{TheOracle I}, \lean{Seam.flock}, \lean{Seam.done},
  \code{Phases.opening : Phase.Def I (I.Stmt × FlockOut I) Unit} stay. No built phase queries the
  oracle (the public-input verifier queries only its own message, \code{PublicInput.lean:222-231}
  at \code{b435631}), so no built proof changes. (That last step is the dossier's inference: every
  type that mentions \code{[TheOracle I]ₒ}, for example \code{PublicInput.lean:242, 246, 261, 263},
  elaborates against whichever \lean{OracleInterface} instance is in scope, so the citation check
  marks it as needing a build to confirm: unverified.) \\
Layer 10
  & the opening phase is the batching component alone (the batching hole, G3, \code{bp:600}):
  schedule \code{V_to_P : E}, verifier = batch, query, compare; error $(J − 1)/\abs{\E}$; knowledge
  state function ``some pooled claim is false of $q$'' before $λ$, ``the batched claim is false of
  $q$'' after, with the query making the verdict; completeness by linearity of \lean{Weight.pair}.
  ``$2/\abs{\E}$ per round'' leaves Layer 10 (those rounds are WHIR's and carry
  $2L_i/\abs{\E} + 2^{ℓ_i − j}ε_i$, Theorem B.7). The title and \code{bp:127} lose ``the opening
  sumcheck'' \\
Layer 11
  & \lean{whirOpen} takes the $J$ weighted claims and \textbf{its level-0 batching challenge is the
  opening phase's $λ$}: the compilation replaces the whole opening phase by \lean{whirOpen}, not
  the query alone. \lean{encode} must be over $\E$ as well as $\K$ (\cref{sec:gen-opening-blueprint-layer11}) \\
Layer 12
  & \code{leanVmIopp = commit' ⟫ bus ⟫ table ⟫ pub ⟫ flock ⟫ whirOpen}, where \code{commit'}
  sends the level-0 codeword (an interleaved $\K$-valued word with position queries) and the front
  \code{bus … flock} is the spine's; then the Merkle compilation, then Fiat--Shamir with grinding
  (\cref{sec:gen-opening-wellformed}) \\
Layer 13
  & unchanged in shape; its hypotheses change with \cref{sec:gen-opening-wellformed} \\
\end{longtable}
}

\paragraph{(iii) End the oracle protocol at the pool.} Drop the opening phase from the oracle
protocol: the composed reduction ends with output statement \code{I.Stmt × FlockOut I} and output
relation \lean{Seam.flock}, and the compilation appends \lean{whirOpen} there. Same transcript as
(ii), same compile theorem; it changes the master theorems' shape (\lean{Seam.done} becomes
\lean{Seam.flock}, five phases instead of six) and leaves the stack's oracle interface unused,
which misleads (an interface that is never queried says nothing). Acceptable; (ii) is preferred
because it keeps the spine's statements and makes the ideal object the specification names.

\paragraph{Recommendation: (ii).} It is the specification's own object (\code{03:112-118};
\code{B:4} ``we commit to a $\K$-valued multilinear, and opening claims are in $\E$''), it matches
ArkLib's \lean{Commitment.Scheme} without adaptation, it makes \lean{verify_iff_compiled} statable
with leanVM's transcript, and it removes a $μ$-round sumcheck from the trusted surface of Layer 10.

% ---------------------------------------------------------------- A.6
\paragraph{Negative results.} Checked and agreeing across the three sources: the moment of $λ$
(after the six ring-switching challenges, before any WHIR message); the power of each claim;
$J = 113$; the batched target and weight; that the verifier evaluates $\mle{W}_λ$ exactly once.
Checked: the spine's \lean{WeightedClaim.Holds} (\code{Seams.lean:123-124}) is
$Σ_w W(w)·q(w) = c$ over the cube with $q$ lifted by \lean{ofK}, the specification's
$⟨W, g⟩ = c$ (\code{03:115}); and \lean{Weight.mle_eq} ties the verifier's evaluator to the cube
values, which is what ``MLE-friendly'' needs. No finding on the pool's translation into weighted
claims.

% ---------------------------------------------------------------- B.1
\subsubsection{WHIR level by level}\label{sec:gen-opening-whir}

Notation: $μ$ the stack's log-size; $k_0 = 6$ the initial fold (\code{whir_config.rs:67}), $k_i = 4$
after (\code{:68}), the residual $n_{\mathrm{res}} ≤ 5$ (\code{:86}); $r$ the number of recursive
levels (\code{level_steps}); $t_i$ the query count of level $i$; $z$ an out-of-domain (OOD) point.
The specification's column is Protocol B.6 (\code{B:100-123}); the Rust column gives the prover and
the (succinct) verifier, both in \code{whir.rs} unless another file is named; the Python column is
\code{py:1006-1087}.

\begin{landscape}
{\footnotesize\setlength{\tabcolsep}{4pt}
\begin{longtable}{@{}>{\raggedright\arraybackslash}p{0.03\linewidth}>{\raggedright\arraybackslash}p{0.31\linewidth}>{\raggedright\arraybackslash}p{0.23\linewidth}>{\raggedright\arraybackslash}p{0.25\linewidth}>{\raggedright\arraybackslash}p{0.12\linewidth}@{}}
\caption{WHIR's message schedule per level, as deployed.}\label{tab:gen-opening-whir}\\
\toprule
\# & Message & Specification, Protocol B.6 & Rust prover / verifier & Python \\
\midrule
\endfirsthead
\toprule
\# & Message & Specification, Protocol B.6 & Rust prover / verifier & Python \\
\midrule
\endhead
\bottomrule
\endfoot
0 & $V→P$ $λ$ (level-0 batching), $J_0 = 113$ claims
  & step 1 at $i = 0$ (\code{B:110})
  & outside WHIR: \code{stack_open.rs:400} / \code{:518} (\cref{sec:gen-opening-lambda})
  & \code{py:1361} \\
1 & $P→V$ $h_1$: round polynomial of $Σ_x q(x)·W_λ(x) = C_λ$ (2 scalars: $c_0, c_2$; $c_1$ derived)
  & step 2, $j = 1$ (\code{B:111})
  & \code{:1324-1326} / \code{:2087-2090}
  & \code{py:1012} \\
2 & 6 × ($V→P$ $r_j$; $P→V$ next round polynomial), the \textbf{lane fold} on the stack's top 6
  coordinates
  & step 2, $j = 1..ℓ_0$ (\code{B:111-112}); the row index is the stack's most significant $ℓ_0$
  coordinates, $f(u, x) := q(x, u)$ (\code{B:380})
  & \code{:1328-1334} / \code{:2100} (\code{replay_fold_rounds}, \code{:1641-1655}); the point is
  rotated by 6 at the end, \code{:2297-2306}
  & \code{py:1017-1024}; rotation \code{py:1078-1079} \\
3 & $P→V$ root of $f^{(1)}$ (2 scalars), the fold $f^{(1)} = \mle{q}(r_1..r_6, ·)$, $\E$-valued,
  interleaved $2^4$ lanes
  & step 3(a) (\code{B:115})
  & \code{:1340-1358} / \code{:2104}
  & \code{py:1034} \\
4 & OOD: $V→P$ $z ∈ \E^{μ−6}$ ($μ − 6$ squeezes); $P→V$ $y$ (1 scalar); $P→V$ intro round
  polynomial for the claim $(\eq(z, ·), y)$ (2 scalars)
  & step 3(b) (\code{B:116}): $z$, $y$; \textbf{no intro message}
  & \code{send_ood}, \code{:1201-1208} / \code{replay_ood}, \code{:1670-1675}; count
  \code{ood_samples[1] = 1} (\code{whir_config.rs:677-684}, test \code{:1032-1034})
  & \code{py:1035-1037} \\
5 & $P→V$ grinding nonce (1 raw scalar); $V$ checks 17 bits
  & \textbf{absent}
  & \code{:1366} / \code{:2117} (\code{grind_check}, \code{fs/transcript.rs:314-323});
  \code{QUERY_GRINDING_BITS = 17} (\code{whir_config.rs:60})
  & \code{py:1039} \\
6 & $V→P$ $t_0$ positions in $[0, 2^{μ−6+ρ})$, by chunking squeezes
  & step 3(c): ``$t_i$ i.i.d.\ columns'' (\code{B:117})
  & \code{sample_queries_ordered}, \code{:1175-1193}: $⌊192/d⌋$ $d$-bit chunks per squeeze,
  duplicates kept / \code{:2122}
  & \code{py:943-951}, \code{py:1041} \\
7 & $V→P$ $λ_0$, the level's batching challenge
  & step 1 at $i = 1$ (\code{B:110}), \textbf{after} step 3(d)'s claims are formed
  & \code{:1373} / \code{:2123}
  & \code{py:1044} \\
8 & $P→V$ Merkle openings of level 0 at the positions (rows of $n_{\mathrm{lanes}}$ words and one
  octopus), \textbf{not absorbed}
  & step 3(c): ``it reads $C^{(i)}[·, x_q]$'' (\code{B:117})
  & \code{hint_merkle}, \code{:1380} / \code{recv_level_rows}, \code{:2125-2135}
  (\code{fs/transcript.rs:263-279})
  & \code{py:1049-1050} (raw paths) \\
9 & $P→V$ intro round polynomial for the query batch (2 scalars), with claim
  \code{enforced_sum_0} $= Σ_q λ_0^q·⟨\mathit{row}_q, \eq(r_{\mathrm{lane}}, ·)⟩$ computed by $V$ from
  the openings
  & step 3(d): the $t_i$ consistency claims $(W_{x_q}, c_q)$ with $c_q = C_{\mathrm{fold}}[x_q]$
  (\code{B:117-118}); \textbf{no intro message}, the batch's $h_1$ is sent after $λ$
  & \code{:1405-1406} / \code{:2140-2144} (\code{induce_sumcheck_enforced_sum},
  \code{whir_induce.rs:203-218})
  & \code{py:1051}, \code{py:1057} \\
10 & $V$ glues: running quad \code{+= λ_0·quad_ood + λ_0²·quad_query}, claim
  \code{t_r += λ_0·y + λ_0²·enforced_sum}
  & step 1 at level 1: $(W, σ_0) := (Σ λ^{t−1} W_t, …)$ with the order residual, OOD, queries
  (\code{B:118}, \code{B:110})
  & \code{glue_pending}, \code{:1131-1154} / \code{batch_level_claims}, \code{:1681-1701}
  & \code{py:1058-1063} \\
11 & level $i ≥ 1$: 4 × ($V→P$ $r$; $P→V$ round polynomial); then as rows 3 to 10 with $t_i$, the
  leaf $3·2^4$ words
  & steps 2--3 (\code{B:111-119})
  & \code{:1420-1425, 1509-1560} / \code{:2191-2194, 2311-2384}
  & \code{py:1017-1063} (loop) \\
12 & last level: after its 4 folds, $P→V$ $f_{\mathrm{final}}$ ($2^{n_{\mathrm{res}}}$ scalars in
  the clear); grind; $V→P$ $t_{\mathrm{last}}$ positions; $V→P$ $λ_{\mathrm{last}}$; openings;
  intro; glue; then $n_{\mathrm{res}}$ × ($V→P$ $r$; $P→V$ round polynomial \textbf{except after
  the last $r$})
  & step 4 (\code{B:120}): $f_{\mathrm{final}}$, queries, $λ$, $ν_L$ rounds, ``$V$ closes by
  evaluating $\mle{f}_{\mathrm{final}}$''; no intro; nothing said about the last message
  & \code{:1431-1476} / \code{:2197-2262}
  & \code{py:1031-1032}, \code{py:1065-1074} \\
13 & terminal check $\mathit{weight} · \mle{f}_{\mathrm{final}}(ρ_{\mathrm{tail}}) = t_r$ with
  $\mathit{weight} = \mle{W}_λ(ρ \text{ rotated}) + Σ_{\text{levels}} β_i·(\text{induced weight of
  the level's queries at } ρ) + Σ_{\text{ood}} β·\eq(z, ρ)$
  & the closing check of step 4 (\code{B:120}), the consistency claims being
  $⟨W_{x_q}, f_{\mathrm{final}}⟩$ (\code{B:120}, Lemma \code{lem:colweight})
  & \code{:2264-2308} / \code{induce_sumcheck_evaluate_at_residual}, \code{whir_induce.rs:228-293}
  & \code{py:1075-1083}, \code{_induced_weight}, \code{py:971-989} \\
\end{longtable}
}
\end{landscape}

\begin{sloppypar}
The two verifiers agree with each other message for message (checked by walking
\code{recursive_verifier_with_basis_succinct} and \code{verify_whir} side by side; the Python is a
transcription of the Rust). Against Protocol B.6 the deployed schedule differs in four ways: the
per-claim intro messages, the grinding, the lane relayout, and the omitted last round message; they are items 6 to 9 of \cref{tab:gen-opening-disagreements}.
\end{sloppypar}

Also: the OOD count is one per level $≥ 1$ and zero at level 0 (\code{whir_config.rs:677-684},
\code{:829-841}; \code{whir.rs:1261-1265} ``L0 takes no OOD sample''), as Annex B (\code{B:105}
``No OOD sample is taken'') and the blueprint (\code{bp:1186-1187}) say.

\paragraph{Element counts} (from \cref{tab:gen-opening-whir}, \cref{sec:gen-opening-code} and
\cref{sec:gen-opening-rounds}).

\begin{center}\small
\begin{tabular}{@{}>{\raggedright\arraybackslash}p{0.42\linewidth}>{\raggedright\arraybackslash}p{0.5\linewidth}@{}}
\toprule
Message or challenge & Elements of $\E$ \\
\midrule
round polynomial of a fold round & 2 scalars, $c_0$ and $c_2$ ($c_1$ derived); its challenge $r$ is
  1 squeeze \\
root of a level & 2 scalars \\
OOD point after the first root & $μ − 6$ squeezes \\
OOD answer $y$ & 1 scalar, one OOD sample per level $≥ 1$, none at level 0 \\
intro round polynomial (OOD claim; query batch) & 2 scalars each \\
grinding nonce & 1 raw scalar at every level \\
query positions & $⌊192/d⌋$ chunks of $d$ bits per squeeze, duplicates kept \\
level batching challenge $λ_i$ & 1 squeeze \\
Merkle openings & none on the stream: a side channel, not absorbed \\
$f_{\mathrm{final}}$ & $2^{n_{\mathrm{res}}}$ scalars \\
round message after the last $r$ & omitted \\
\bottomrule
\end{tabular}
\end{center}

The ladder of folds is $[6, 4, 4, …]$ until at most 5 variables remain: 2 levels and a residual of
5 at $μ = 15$, 6 levels and a residual of 2 at $μ = 28$ (\cref{sec:gen-opening-code}). The query
counts per level are those of \code{py:910} (at $μ = 15$ and rate $1/2$: 223, then 55).

% ---------------------------------------------------------------- B.2
\subsubsection{The code}\label{sec:gen-opening-code}

\begin{sloppypar}
\begin{itemize}
\item \textbf{Field and domain.} Reed--Solomon over $\E$ with the evaluation domain an
  $\F_2$-subspace of $\K$ (\code{B:54}, \code{B:90}); the level-0 word is $\K$-valued because the
  message is (\code{B:306}; the codeword buffer is \code{ArenaVec<F64>}, \code{whir.rs:293},
  \code{:364}); deeper levels are $\E$-valued (\code{ArenaVec<F192>}, \code{:397}) encoded with $\K$
  twiddles by the mixed product (\code{whir.rs:17-19}, \code{whir_ntt_ext.rs}).
\item \textbf{Basis.} The novel polynomial basis of Lin--Chung--Han on the standard basis
  $e_c = x^c$ of $\K$ (\code{B:278}, \code{B:301-306}): \code{whir_induce.rs:19-21} ``Standard basis
  only (\code{v_i = x^i = F64(1 << i)})'', \code{eval_sk_at_vks} (\code{:32-52}) and
  \code{normalized_sks_at} (\code{:56-67}, $\hat{W}_i(x) = s_i(x)/s_i(v_i)$, the specification's
  $\hat{W}_i := W_i/W_i(e_i)$, \code{B:299}); the Python \code{_subspace_roots} (\code{py:962-968}).
  A position $q ∈ [0, 2^{κ+ρ})$ is the domain point \code{F64(q)} (\code{whir_induce.rs:163},
  \code{py:983}), the specification's $x_v = Σ_c v_c e_c$ (\code{B:357}). The message is read
  directly as novel-basis coefficients (\code{B:306}; the encoder is the additive NTT,
  \code{ntt::AdditiveNttF64::standard}, \code{whir.rs:372-373}).
\item \textbf{The column weight} the verifier evaluates: $∏_k (1 + p_k·(1 + \hat{W}_k(q)))$
  (\code{whir_induce.rs:220-224}, \code{py:971-977}), which is Lemma \code{lem:colweight}'s
  $∏((1 + r_i) + r_i \hat{W}_i(x))$ (\code{B:312}) in characteristic two. Agrees.
\item \textbf{Rates.} Level 0: $ρ_0 = 2^{−ρ}$ with $ρ ∈ [1, 4]$ chosen by the prover and announced
  (\code{whir_config.rs:41-55}, \code{mod.rs:123, 149, 167}); then the rate exponent grows by
  $k_0 − 3 = 3$ after the initial fold and by $k_i − 1 = 3$ after each later fold
  (\code{whir_config.rs:73-78, 298-311}; \code{py:929}). Block lengths $2^{κ_i + ρ_i}$, all
  $≤ 2^{28}$, within $κ + ρ ≤ 64$ (\code{B:90}).
\item \textbf{The ladder.} Folds $[6, 4, 4, …]$ until at most 5 variables remain
  (\code{derive_ladder}, \code{whir_config.rs:260-296}; \code{py:924-932}): 2 levels and a residual
  of 5 at $μ = 15$, 6 levels and a residual of 2 at $μ = 28$.
\end{itemize}
\end{sloppypar}

% ---------------------------------------------------------------- B.3
\subsubsection{The parameters}\label{sec:gen-opening-parameters}

\begin{sloppypar}
\code{SECURITY_BITS = 128}, \code{LOG_INV_RATE_0 = 1}, \code{MIN/MAX_LOG_INV_RATE = 1/4},
\code{QUERY_GRINDING_BITS = 17}, \code{INITIAL_FOLDING_FACTOR = 6},
\code{SUBSEQUENT_FOLDING_FACTOR = 4}, \code{RS_DOMAIN_INITIAL_REDUCTION_FACTOR = 3},
\code{RS_DOMAIN_SUBSEQUENT_REDUCTION_FACTOR = 1}, \code{RESIDUAL_MAX_LOG = 5}
(\code{whir_config.rs:38-86}); \code{MIN_MU = 15}, \code{MAX_MU = 28} (\code{pcs.rs:49-51}); the
Python's copies \code{py:900-908}. The query counts per level come from the search of
\code{whir_config.rs:639-707} (\cref{sec:gen-opening-deployed-claim}) and are tabulated in
\code{py:910} for every $(ρ, μ)$; the scratch port reproduces the table
(\cref{sec:gen-opening-probe-params}).

The blueprint's parameter block (\code{bp:1157-1165}) lists the six constants and
\code{ladder … -- fold, rate exponent, queries per level} with \lean{ladder_queries_eq} against
\code{py:910}: right, with two gaps. (a) The analysis parameter $η_i$ (or $m_i$) per level is not a
field of \lean{Level}, yet \code{whirError params} (\code{bp:1175}) is a function of
$γ_i = 1 − \sqrt{ρ_i} − η_i$ and $L_i = 1/(2η_i\sqrt{ρ_i})$; a Lean \lean{whirError} needs the $η_i$
as data (any admissible value proves the theorem; the Rust's are the floating-point optimum,
\code{whir_config.rs:624-633}), and acceptance of ``128 bits'' needs the specific ones. (b) The OOD
count (1 from level 1) and the grinding bits per level are parameters too
(\cref{f:opening-level-eta}).
\end{sloppypar}

% ---------------------------------------------------------------- B.4
\subsubsection{The Merkle trees and the encodings}\label{sec:gen-opening-merkle}

\begin{sloppypar}
\begin{itemize}
\item \textbf{Tree.} Binary, $2^k$ leaves, flat layout (\code{pcs/merkle.rs:11-18}); leaf hash =
  BLAKE2s-256 of the leaf bytes (\code{fs/merkle.rs:38-41}), node = BLAKE2s-256 of the 64 bytes
  $\mathit{left} ‖ \mathit{right}$ (\code{:44-51}); \textbf{no domain separation} between leaves and
  nodes. Safe because the verifier fixes the height (\code{open}, \code{fs/merkle.rs:174}:
  \code{height = num_leaves.trailing_zeros()}, \code{num_leaves} a function of the announced sizes)
  and the leaf width (\code{row_words}, \code{leaf_words}, \code{:139-148}); a node's 64-byte
  preimage can only be presented as a leaf of a tree whose height the verifier does not accept.
  (The citation check gives the more direct reason, the length: a leaf is BLAKE2s-256 of 512 bytes at
  level 0 or 384 bytes deeper, eight or six chained compressions, the committer relying on the
  chaining value after the leading zero blocks, \code{pcs/merkle.rs:147-155}; a node is one 64-byte
  compression.)
\item \textbf{Leaf images.} Level 0: $2^6$ words of 8 bytes (512 bytes) per position,
  lane-\textbf{descending}, the absent lanes a zero prefix supplied by the verifier
  (\code{whir.rs:331-345}, \code{pcs/merkle.rs:144-192}, \code{fs/merkle.rs:53-58}); the words are
  little-endian \code{u64}s (\code{fs/merkle.rs:60-67}). Deeper levels: $2^4$ elements of $\E$ as
  $3·2^4 = 48$ words, 384 bytes (\code{whir.rs:461-467}, \code{ext_row_words}, \code{:1211-1213};
  \code{py:1049}, \code{_ext_row}, \code{py:938-940}). The comment \code{whir.rs:393-395} says ``one
  Merkle leaf of \code{num_interleaved * 16} bytes''; the code hashes
  \code{num_interleaved * size_of::<F192>() = 24·num_interleaved} bytes (\code{:463}); the comment is
  stale (\cref{f:opening-stale-leaf-comment}).
\item \textbf{Digests on the stream.} A 32-byte digest travels as two scalars, each two
  little-endian words in the low limbs and a zero third limb (\code{hash_to_scalars},
  \code{fs/merkle.rs:14-20}; \code{py:216-219}); reading rejects a nonzero third limb
  (\code{scalars_to_hash}, \code{:26-36}; \code{py:222-224}).
\item \textbf{Wire format.} One phase per level: the rows at the \emph{distinct, sorted} positions
  and one ``octopus'' of siblings (\code{PrunedMerklePaths}, \code{fs/merkle.rs:78-92}, \code{prune},
  \code{:101-136}); the verifier rebuilds every path, refuses a missing or surplus sibling, a wrong
  row count or width, an out-of-range position, a root mismatch (\code{open}, \code{:163-235}). The
  raw form (\code{RawMerklePath}, \code{:246-264}: full leaf image plus full sibling path, one per
  query, duplicates repeated) is what an accepting Rust verifier \emph{emits}
  (\code{into_raw_proof}, \code{fs/transcript.rs:198-203}; \code{mod.rs:702-704, 777}) and what the
  Python and the recursion guest consume (\code{py:392-404}; the test harness
  \file{crates/lean_vm/tests/verifiers/python_verifier.rs:50-80} writes it out).
\end{itemize}
\end{sloppypar}

% ---------------------------------------------------------------- B.5
\subsubsection{What Annex B's soundness theorem says, and on what it rests}\label{sec:gen-opening-thm-rbr}

Theorem B.7 (\code{B:139-152}, \code{thm:rbr}): the scheme is $L_0$-list binding (Definition B.4,
\code{B:70-72}) and its opening has \textbf{round-by-round soundness} (Definition B.5,
\code{B:74-84}) for the relation $R_{\mathrm{open}}$ (\code{B:133-136}), with the per-message errors
of the table \code{B:141-150}. The model is the interactive oracle proof: the oracles $C^{(i)}$ are
read at queried columns, the challenges are uniform, the prover unbounded; the extractor is not
named (soundness for a language). Inputs: the Johnson bound, proved in the annex
(\code{B:391-414}); mutual correlated agreement up to the Johnson bound, Theorem B.9
(\code{B:176-185}), quoted from [BCHKS25] Theorem 4.6 and external; Lemma B.10 (\code{B:189-198})
turns it into ``folding preserves lists'' with the $2^{ℓ−1}$ union of rows; Lemma B.11
(\code{B:200-206}) is the out-of-domain separation. Fiat--Shamir is one sentence (\code{B:84}).

% ---------------------------------------------------------------- B.6
\subsubsection{The blueprint's Layer 11 set against it}\label{sec:gen-opening-blueprint-layer11}

{\footnotesize\setlength{\tabcolsep}{4pt}
\begin{longtable}{@{}>{\raggedright\arraybackslash}p{0.31\linewidth}>{\raggedright\arraybackslash}p{0.27\linewidth}>{\raggedright\arraybackslash}p{0.36\linewidth}@{}}
\caption{The blueprint's Layer 11 (\code{bp:1150-1193}) against the deployed WHIR.}\label{tab:gen-opening-layer11}\\
\toprule
Blueprint & Ground truth & Judgement \\
\midrule
\endfirsthead
\toprule
Blueprint & Ground truth & Judgement \\
\midrule
\endhead
\bottomrule
\endfoot
\code{whirOpen (params) : OracleReduction … (StmtIn := Fin J → WeightedClaim μ) (OStmtIn := codeword oracles)},
  ``stated exactly as Protocol B.1: batch, $ℓ_i$ sumcheck rounds, commit, one out-of-domain sample
  from level 1 on, $t_i$ queries, and the final plaintext level''
  & Protocol B.6 (numbered on the section counter, \code{theorems.tex:4, 12}) with the four deployed differences of
  \cref{sec:gen-opening-whir}
  & the input type is right ($J$ weighted claims; under (ii) its level-0 batch is the opening
  phase's $λ$, \cref{sec:gen-opening-alternatives}). The relayout, the intro messages and the
  omitted last message must be pinned as transcribed content (the blueprint's Category B) for
  \lean{verify} to accept Rust proofs; grinding is Layer 12's. ``Protocol B.1'', ``Theorem B.2'',
  ``Definition B.1'' (\code{bp:1188}; the citation check corrected the dossier's \code{bp:1187}), ``Lemma B.7'' (\code{bp:1170}), ``Annex B.3''
  (\code{bp:278}), ``(Annex B.4)'' (\code{bp:1172}) are not the pinned numbers (Protocol B.6,
  Theorem B.7, Definition B.4, Lemma B.14, §B.5, §B.3); cite labels \\
\code{whirOpen_rbrSoundness (mca : McaJohnson) : … (whirError params) -- Theorem B.2, per message};
  ``Its soundness is \emph{round-by-round soundness for the list relation} \code{relopen}, not
  knowledge soundness''
  & Theorem B.7 as read in \cref{sec:gen-opening-thm-rbr}
  & right in kind; \lean{whirError} needs the $η_i$ (\cref{sec:gen-opening-parameters}); the errors
  are $(J_i − 1)L_i/\abs{\E}$, $2L_i/\abs{\E} + 2^{ℓ_i−j}ε_i$, $C(L_i, 2)\,μ_i/\abs{\E}$,
  $(1 − γ_i)^{t_i}$, $t_{r−1}/\abs{\E}$, $2/\abs{\E}$; Layer 10's ``$2/\abs{\E}$ per round'' is
  the tail's value, wrong for the fold rounds (\cref{sec:gen-opening-blueprint}) \\
\code{encode (κ R : ℕ) : Column κ → (Fin (2^(κ+R)) → K) -- additive NTT on the K basis}
  & the encoder is $\mathrm{Enc} : \E^{\mathrm{cube}} → \E^{\mathrm{dom}}$, $\K$-valued on
  $\K$-valued messages (\code{B:54, 306}); every level past 0 encodes an $\E$-valued message
  (\code{whir.rs:419-485})
  & \textbf{too narrow}: \lean{encode} over $\K$ serves level 0 only; the deeper levels need the
  same encoder over $\E$ (or over any $\K$-algebra), and the $\K$-valuedness of level 0 is a lemma
  about it (\cref{f:opening-whiropen-batch-encode}) \\
\code{theorem encode_column_weight (x) : ∃ W_x : Weight κ, ∀ f, encode κ R f x = ⟨W_x, f⟩ -- Lemma B.7}
  & Lemma B.14 (\code{B:309-315}): $W_x(u) = ∏_i \hat{W}_i(x)^{u_i}$ and
  $\mle{W}_x(r) = ∏_i ((1 + r_i) + r_i \hat{W}_i(x))$, which the verifier computes
  (\code{whir_induce.rs:220-224})
  & \textbf{near-vacuous as stated}: every linear map $\K^{2^κ} → \K$ is an inner product with some
  cube table, and \lean{Weight} is inhabited by \code{⟨row, evalMle row, rfl⟩}, so the existential
  is true of any linear \lean{encode} and says nothing about the novel basis. What Layer 11 needs
  is the \emph{definition} \code{columnWeight x : Weight κ} with the closed-form \lean{mle} and the
  theorem \code{encode κ R f x = (columnWeight x).pair f}, since \lean{verify} evaluates that
  closed form (\cref{f:opening-encode-column-weight}) \\
\code{merkleRoot : List (Vector K leafWords) → Digest},
  \code{merkleVerify : Digest → ℕ → Vector K leafWords → Path → Bool}, \code{blake2sBytes}
  & \cref{sec:gen-opening-merkle}
  & \lean{merkleVerify} is the raw path (\code{RawMerklePath::root}, \code{fs/merkle.rs:252-264}),
  not the wire's pruned octopus; which one \code{Proof.merkle : List MerklePaths}
  (\code{bp:1209-1211}) holds decides which proofs \lean{verify} accepts
  (\cref{sec:gen-opening-proof}). The leaf image order (lane-descending, zero prefix) and the two
  leaf widths are not stated \\
Tests: \lean{encode} against CompPoly's additive NTT at $κ = 3$; \lean{merkleVerify} on four
  leaves; honest \lean{whirOpen} at $μ = 4$, rate $1/2$; the parameter tables against
  \code{py:910}
  &
  & the honest run at $μ = 4$ cannot use the production ladder (\code{derive_ladder} needs
  $μ > 6$ and two levels, \code{whir_config.rs:266-267, 291-293}; the Rust's tests fall back to
  \code{default_config}, \code{:196-224}; the fallback selector is \code{test_configs_for},
  \code{:232-244}, as the citation check adds); say which shape the toy uses \\
\end{longtable}
}

% ---------------------------------------------------------------- C
\subsubsection{Fiat--Shamir: the chain}\label{sec:gen-opening-chain}

§8.4 of the specification is \code{TODO} (\code{08:44-46}); §8.5's ``Setup'' says only ``The
Fiat-Shamir initial state (§8.4) is seeded by the public input, the bytecode and Flock's BLAKE2s
R1CS matrices'' (\code{08:55}). \textbf{The Rust is the only source}; the Python is a transcription
of it and is checked against it below (every row agrees unless marked).

{\footnotesize\setlength{\tabcolsep}{4pt}
\begin{longtable}{@{}>{\raggedright\arraybackslash}p{0.17\linewidth}>{\raggedright\arraybackslash}p{0.62\linewidth}>{\raggedright\arraybackslash}p{0.15\linewidth}@{}}
\caption{The Fiat--Shamir chain, as deployed.}\label{tab:gen-opening-chain}\\
\toprule
Element & Rust (\code{fs/lib.rs} unless named) & Python \\
\midrule
\endfirsthead
\toprule
Element & Rust (\code{fs/lib.rs} unless named) & Python \\
\midrule
\endhead
\bottomrule
\endfoot
The primitive \code{compress(a, b)} on two 256-bit halves
  & \code{:18-24}: $\mathrm{BLAKE2s}(a ‖ b)$ over 64 bytes, the words little-endian, that is the RFC
  compression $F(\mathrm{PARAM\_IV}, a ‖ b, t = 64, f_0 = 1)$ from the unkeyed parameter state
  (\code{primitives/src/hash.rs:83-87, 209-224}); the chain state is the \emph{first half of the
  message}, never a BLAKE2s chaining value
  & \code{py:341-343} \code{compress(left, right) = blake2s(le bytes of the 8 words)} \\
Tags, lane 3 of the right half
  & \code{:36-39}: observe \code{1}, squeeze \code{2}, grinding base \code{3}, grinding nonce
  \code{4}; the seeding block has no tag, ``its position is its tag'' (\code{:31-35})
  & \code{py:335-338} \\
Seed
  & $\mathit{cv}_0 = \mathrm{compress}(\mathit{iv}, \mathit{public\_input})$ (\code{:69-73});
  \code{iv = BLAKE2s("leanvm" ‖ len(R1CS_DIGEST) as u64 LE ‖ R1CS_DIGEST ‖ bytecode_hash)} as four
  words (\code{mod.rs:82-93}, \code{:712}); the public input as its four low words, the third limbs
  dropped (\code{mod.rs:99-106})
  & \code{py:1366-1369}, \code{py:349} \\
Observe $x ∈ \E$
  & \code{cv ← compress(cv, (x.c0, x.c1, x.c2, 1))} (\code{:85-87})
  & \code{py:353-354} \\
Squeeze
  & \code{out ← compress(cv, (0, 0, 0, 2)); cv ← out}; challenge \code{= (out_0, out_1, out_2)}
  (\code{:95-99}): the challenge is 192 of the 256 state bits
  & \code{py:356-358} \\
Grinding base
  & \code{compress(cv, (0, 0, 0, 3))}, not written to the state (\code{:108-110})
  & \code{py:380} \\
Nonce bound
  & \code{cv ← compress(cv, (nonce.c0, nonce.c1, nonce.c2, 4))} (\code{:118-120}), on both sides,
  \textbf{also when the check fails} (\code{:165-174}; the status file's finding F13)
  & \code{py:382} \\
Grinding predicate
  & the low \code{bits} bits of \code{compress(base, (nonce, 4))[0]} are zero; \code{bits = 0}
  accepts only \code{nonce = 0} (\code{:47-51}, \code{:165-174})
  & \code{py:381} \\
The honest grind
  & the smallest \code{u64} nonce, unbounded search (\code{:126-155})
  & not the verifier's \\
\end{longtable}
}

The blueprint's row \emph{Fiat--Shamir} (\code{bp:325}) states the tags, the seed and ``one squeeze
yields one $\E$ challenge (three low words)'' correctly. It does not say that the step is
$\text{BLAKE2s-256}(\mathit{cv} ‖ \mathit{block})$ with $\mathit{cv}$ in the \emph{message} and the
RFC chaining value fixed at the parameter state; Layer 12's \code{FsState.observe … -- lane 3 = 1}
(\code{bp:1204}) and the leanISA row ``\code{compress} (Layer 1): the compression the transcript,
Merkle hashing and the Flock circuit share'' (\code{bp:216}) invite
$F(\mathit{cv}, x ‖ \mathit{tag}, …)$ with $\mathit{cv}$ as chaining value, which is a different
function. It is transcribed content (the blueprint's Category B); the transcription must say
\code{hash64 (cv ++ x ++ tag) = compress iv64 (cv ++ x ++ tag) 64 true false} in the terms of
leanISA's \lean{compress} (\file{LeanerVM/Semantics/Blake2s.lean:105} at \code{b435631})
(\cref{f:opening-chain-step}). The citation check qualifies both halves. \code{bp:216} also says ``the
byte hasher … is the RFC 7693 wrapper over \code{compress} with the last-block flag only'', which the
dossier does not mention, so ``invite'' is an inference; and the proposed transcription is imprecise:
leanISA's \lean{compress} takes \code{(h : Vector UInt32 8) (m : Vector UInt32 16) (t : UInt64) (f0 f1 : UInt32)},
the flags being \code{0xFFFFFFFF} when set, so \code{true false} does not typecheck, and \code{iv64}
must be the parameter-folded IV (\code{PARAM_IV}, \code{primitives/src/hash.rs:64-81}), not the bare
IV. The source's own field comment calls \code{cv} ``The 256-bit chaining value: a Merkle-Damgård hash
of the transcript so far'' (\code{fs/lib.rs:58}), in the Merkle--Damgård sense, which the check
names as the likely origin of the ambiguity.

\subsubsection{The proof object and the stream order}\label{sec:gen-opening-proof}

\begin{sloppypar}
\code{Proof<M = PrunedMerklePaths> { stream: Vec<F192>, merkle: Vec<M> }}
(\code{fs/transcript.rs:9-12}); \code{RawProof = Proof<RawMerklePath>} (\code{:19}) is what the Rust
verifier emits after accepting (\code{:198-203}) and what the Python reads (\code{py:322-330}: a byte
file of 24-byte scalars and a byte file of leaf words followed by sibling digests, no lengths). The
blueprint's \lean{Proof} (\code{bp:1209-1211}: \code{stream : List E}, \code{merkle : List MerklePaths})
cites \code{transcript.rs:9-19}, which defines both; \code{merkleVerify … Path} (\code{bp:1182}) is
the raw shape; the fixture ``a proof dumped from the pinned Rust prover'' (\code{bp:1237-1238}): the
prover emits the pruned form (\code{ps.into_proof()}, \code{:138-143}), the verifier the raw one.
Which one \lean{verify} takes is unspecified, and they are not equivalent as sets of accepted byte
strings:
\begin{enumerate}[label=(\alph*)]
\item the raw form repeats a row for a duplicated position and carries every sibling, the pruned
  form dedups and shares (\code{fs/merkle.rs:69-76, 101-136});
\item at level 0 the pruned form carries $n_{\mathrm{lanes}}$ words per row and the verifier
  supplies the zero prefix (\code{fs/merkle.rs:53-58}; \code{whir.rs:2130-2131}), while the raw form
  carries the full 64-word image (\code{py:1049}: \code{lanes if level == 0 else 3 * lanes}) and
  \textbf{the Python does not check that the prefix is zero}: run standalone on an adversarial raw
  proof it accepts a commitment whose absent lanes are nonzero (soundness-neutral, check M2 of
  \cref{tab:gen-opening-checks}; a divergence in accepted sets nonetheless). The citation check
  agrees by reading, not by running: the Python never checks that a word is zero and has no notion of
  $n_{\mathrm{lanes}}$, so such a proof passes its Merkle check; that the rest of its verification
  then passes is the dossier's argument, plausible and unverified by execution.
\end{enumerate}
A choice must be made and named: the pruned wire format is ``what the Rust prover emits''; the raw
one is what the recursion guest verifies (\cref{f:opening-merkle-format}).
\end{sloppypar}

The stream, in order (\code{mod.rs:711-769}; the phases' internals are in the sibling dossiers
and in \cref{sec:gen-bus-transcript,sec:gen-table-pub-transcript,sec:gen-flock-ring-transcript}):

{\footnotesize\setlength{\tabcolsep}{4pt}
\begin{longtable}{@{}>{\raggedright\arraybackslash}p{0.03\linewidth}>{\raggedright\arraybackslash}p{0.35\linewidth}>{\raggedright\arraybackslash}p{0.27\linewidth}>{\raggedright\arraybackslash}p{0.29\linewidth}@{}}
\caption{The proof stream, as deployed.}\label{tab:gen-opening-stream}\\
\toprule
\# & Scalars read ($P→V$, each absorbed as read) & Squeezes ($V→P$) & Source \\
\midrule
\endfirsthead
\toprule
\# & Scalars read ($P→V$, each absorbed as read) & Squeezes ($V→P$) & Source \\
\midrule
\endhead
\bottomrule
\endfoot
1 & seed & — & \code{mod.rs:712} \\
2 & 8: $κ_{\mathrm{mem}}, τ_0 … τ_5, ρ$ & — & \code{mod.rs:144-149}; \code{py:1372} \\
3 & 2: the commitment root & — & \code{mod.rs:714}, \code{pcs.rs:136-138}; \code{py:1382} \\
4 & bus: 2 roots; per GKR layer the round polynomials and children
  & $α$ (4), $β$, per layer the combiners and the two combination challenges; then 5 boundary
  evaluations read
  & the bus dossier, \file{gt-bus.md}, sections A.3--A.4 (\cref{sec:gen-bus-phase}) \\
5 & table sumcheck: $τ_{\max}$ round polynomials (3 scalars each), then 104 column values
  & $ξ$, then $τ_{\max}$ challenges
  & \file{gt-table-pub.md}, section A.1 \\
6 & 2: $c_0, c_1$ & $r_{\mathrm{pi}}$ first & \code{mod.rs:745-749}; \code{py:1398-1399} \\
7 & Flock: zerocheck and lincheck messages & their challenges & \file{gt-flock-ring.md}, section 2.2 \\
8 & — & 6 ring-switching challenges & \code{stack_open.rs:394} / \code{:513}; \code{py:1343} \\
9 & — & $λ$ & \code{stack_open.rs:400} / \code{:518}; \code{py:1361} \\
10 & WHIR: as rows 1 to 13 of \cref{tab:gen-opening-whir}; Merkle phases on the side channel, one
  per level
  & as \cref{tab:gen-opening-whir}
  & \code{whir.rs:2087-2308}; \code{py:1006-1087} \\
11 & \code{vs.finish()}: the stream and the phase list fully consumed & —
  & \code{mod.rs:769}, \code{fs/transcript.rs:213-219}; \code{py:1414} \\
\end{longtable}
}

What is absorbed before each challenge is therefore: everything read from the stream before it
(including the grinding nonces, with tag 4) and every earlier squeeze (which ratchets the state),
and \textbf{nothing else}: the Merkle openings are never absorbed (\code{fs/transcript.rs:223-224},
the doc comment of \code{hint_merkle} at \code{:225-227}; the citation check corrected the dossier's
\code{:225-228}),
the derived coefficient of a round polynomial is never absorbed (\code{:56-62} ``Binding it would
add nothing''), the table sumcheck's target is never sent (the table-sumcheck dossier,
\file{gt-table-pub.md}), the program enters through its hash and the public input through its four
low words at the seed. A challenge is a function of (statement, sizes, all scalars so far, the
count of squeezes so far).

\subsubsection{The encoding of round polynomials}\label{sec:gen-opening-rounds}

\code{add_round_poly(coeffs, eq)} sends every coefficient but one, ``constant first''
(\code{fs/transcript.rs:56-71}): with \code{eq = false} it drops $c_1$, with \code{eq = true} it drops
$c_0$. \code{next_round_poly(n, claim, eq)} (\code{:289-309}) reads the others in index order,
derives the missing one ($c_1 = \mathit{claim} + Σ_{i≥2} c_i$ for a plain round,
$h(0) + h(1) = \mathit{claim}$ in characteristic two; $c_0 = \mathit{claim} + r·Σ_{i≥1} c_i$ for a
round whose equality factor $r$ the verifier holds) and binds the read ones. The Python
\code{sumcheck_round_poly} (\code{py:406-413}) is the same. WHIR's rounds are plain and quadratic:
two scalars $(c_0, c_2)$ per round (\code{whir.rs:509-522}; \code{py:1012}). The blueprint's row
\emph{Sumcheck messages} (\code{bp:315}) and acceptance test 8 (\code{bp:1287-1289}: ``\lean{decode}
reconstructs $c_1$'') state this for the table sumcheck; \code{RoundPoly.decode (d) (claim) (eq? : Option E)}
(\code{bp:1212}) has the right signature. Consequence for the oracle protocol: the round identity
is not a check in the deployed verifier (check B15 of \cref{tab:gen-bus-checks}; the table-sumcheck
dossier's ``three and none on the wire'', \cref{f:table-pub-round-message}); the oracle protocol's
explicit check is what \lean{decode} makes true by construction.

\subsubsection{The blueprint's Layer 12 set against it}\label{sec:gen-opening-blueprint-layer12}

{\footnotesize\setlength{\tabcolsep}{4pt}
\begin{longtable}{@{}>{\raggedright\arraybackslash}p{0.34\linewidth}>{\raggedright\arraybackslash}p{0.6\linewidth}@{}}
\caption{The blueprint's Layer 12 against the deployed chain and proof object.}\label{tab:gen-opening-layer12}\\
\toprule
Blueprint & Judgement \\
\midrule
\endfirsthead
\toprule
Blueprint & Judgement \\
\midrule
\endhead
\bottomrule
\endfoot
\lean{FsState} with \lean{seed}, \lean{observe}, \lean{sample}, \lean{grind} (\code{bp:1202-1206})
  & right shape; \code{grind (bits) (nonce) : FsState → Bool × FsState} must bind the nonce whether
  or not the check passes (\code{fs/lib.rs:165-174}), and \lean{compress} must be the 64-byte hash
  (\cref{sec:gen-opening-chain}) \\
\lean{Proof} (\code{bp:1209-1211})
  & which Merkle format (\cref{sec:gen-opening-proof}) \\
\lean{RoundPoly.decode} (\code{bp:1212})
  & right \\
\code{verify … phase by phase in the stream order of the conventions} (\code{bp:1230})
  & the conventions' rows give neither the order inside WHIR (\cref{tab:gen-opening-whir}) nor the
  Merkle side channel \\
\lean{verify_iff_compiled} (\code{bp:1218-1221})
  & \cref{sec:gen-opening-iff-compiled} \\
Tests: six mutations, ``a flipped stream scalar in each phase, a wrong Merkle sibling''
  (\code{bp:1238-1239})
  & add: a nonce that fails the grind; a nonzero third limb in a root; a nonzero upper limb in a
  size; an extra trailing scalar; an extra sibling; a level-0 row of the wrong width (the bus
  dossier asks for three of these, \cref{f:bus-canonical-checks}) \\
\end{longtable}
}
